\chapter{From Strong to Weak: Integration by Parts}
\label{ch:weak}

% local vector-field macros not in simbook.sty (trial/test fields)
\providecommand{\vu}{\mathbf{u}}
\providecommand{\vv}{\mathbf{v}}
\providecommand{\vp}{\mathbf{p}}
\providecommand{\vz}{\mathbf{z}}

\chapterorient{The previous chapter handed us five partial differential
equations, each written in \emph{strong form}: a differential operator set
equal to a source, holding at every point of the domain. A computer cannot
discretize that equation directly---it asks too much of the unknown, and it
asks it pointwise. This chapter performs the single move that makes the finite
element method possible: multiply by a test function, integrate over the
domain, and integrate by parts. The derivative is ``spread around,'' the
smoothness demand is halved, and the equation reorganizes into a
\emph{bilinear form} on the left and a \emph{linear functional} on the right.
We do this slowly in one dimension, carefully in three, and then for the curl
operator the vector problems need. The payoff is a single template,
$a(u,v)=(Q\,\diffop u,\diffop v)$, that covers every term every Palace driver
assembles.}

\section{Why weak forms}

A partial differential equation in strong form is a demanding object. Take the
electrostatic equation from \cref{ch:reductions-pde},
\begin{equation}
  -\Div(\varepsilon\,\Grad u) = f \quad\text{in } \dom,
  \label{eq:strong-es}
\end{equation}
and read it literally. It asserts that a particular second-order differential
operator, applied to the unknown potential $u$, equals the source $f$ at
\emph{every single point} of the domain $\dom$. For that assertion even to make
sense, $u$ must be twice differentiable: we have to be able to take its
gradient, multiply by the material coefficient $\varepsilon$, take the
divergence of the result, and have a genuine function come out. If the
permittivity $\varepsilon$ jumps across a material interface---as it does in
every real capacitor, where a dielectric meets a vacuum---then $\varepsilon\Grad
u$ is discontinuous, its divergence is not an ordinary function at the
interface, and \cref{eq:strong-es} fails to hold there in any classical sense.
The strong form has painted us into a corner: the physically correct solution
is not smooth enough to satisfy the equation we wrote down.

\begin{physicsbox}
Across the boundary between two dielectrics, the physics does \emph{not} ask
the field to be smooth. It asks for two interface conditions: the tangential
electric field is continuous, and the \emph{normal} component of the flux
$\varepsilon\Grad u$ is continuous---while the normal derivative of $u$ itself
jumps. A formulation that bakes in twice-differentiability cannot represent this
kink. The weak form, by contrast, treats the interface conditions as
\emph{consequences} rather than assumptions, and asks only that $u$ have one
finite derivative. It is the formulation that matches the physics.
\end{physicsbox}

The weak form escapes the corner with a classical trick. Rather than demand that
a differential equation hold pointwise, we test it against a family of probe
functions and demand only that the \emph{averages} match. Concretely: pick a
\emph{test function} $v$, multiply the equation by it, and integrate over the
domain. If \cref{eq:strong-es} holds pointwise then certainly
\begin{equation}
  \int_\dom \bigl(-\Div(\varepsilon\Grad u)\bigr)\,v \dV
  = \int_\dom f\,v \dV
  \label{eq:weighted-residual}
\end{equation}
holds for every $v$. So far we have only lost information---we replaced a
pointwise statement with a family of integral statements. The decisive step
comes next: integrate the left side by parts. One derivative moves off $u$ and
onto $v$. The second derivative on $u$ disappears, the requirement that $u$ be
twice differentiable evaporates, and what remains asks $u$ to have only
\emph{one} derivative---a demand the discontinuous-flux solution can meet.

\begin{figure}[t]
\centering
\begin{tikzpicture}[node distance=10mm]
  \node[stage, text width=30mm, align=center] (strong)
    {\textbf{strong form}\\[2pt]\footnotesize $-\Div(\varepsilon\Grad u)=f$\\[1pt]\scriptsize needs $u\in$ ``$C^2$''};
  \node[opbox, right=14mm of strong, text width=24mm] (test)
    {multiply by $v$,\\ integrate};
  \node[opbox, right=12mm of test, text width=26mm] (ibp)
    {integrate\\ by parts};
  \node[product, right=14mm of ibp, text width=30mm, align=center] (weak)
    {\textbf{weak form}\\[2pt]\footnotesize $a(u,v)=\ell(v)$\\[1pt]\scriptsize needs $u\in$ ``$C^1$''};
  \draw[flow] (strong) -- (test);
  \draw[flow] (test) -- (ibp);
  \draw[flow] (ibp) -- (weak);
  % derivative-count annotation underneath
  \node[font=\scriptsize\itshape, simRust, below=7mm of test.south] (d2) {2 derivatives on $u$};
  \node[font=\scriptsize\itshape, simGreen, below=7mm of weak.south] (d1) {1 on $u$, 1 on $v$};
  \draw[refedge] (test.south) -- (d2);
  \draw[refedge] (weak.south) -- (d1);
\end{tikzpicture}
\caption[Strong to weak: spreading the derivative]{The strong-to-weak pipeline.
Multiplying by a test function $v$ and integrating turns a pointwise equation
into an averaged one; integrating by parts then moves one derivative off the
unknown $u$ and onto the test function $v$. Two derivatives concentrated on $u$
become one derivative on each---the smoothness demand on the unknown is halved,
which is exactly what lets the solution have the kinks the physics requires.}
\label{fig:strongweak}
\end{figure}

There is a second reason to prefer the weak form, and it is not merely
technical: it is the form in which the physics was \emph{originally} stated.
Long before anyone wrote Poisson's equation as a differential equation,
physicists characterized equilibrium by a variational principle. A soap film
settles into the shape of least area; a hanging chain into the shape of least
potential energy; an electrostatic field into the configuration that minimizes
the stored energy $\tfrac12\int_\dom \varepsilon\,\abs{\Grad u}^2\dV$ subject to
the imposed voltages. The strong-form PDE is what you get by differentiating
that energy and setting the result to zero---the Euler--Lagrange equation of the
variational problem. The weak form sits one step \emph{before} the
differentiation: it is the statement that the energy is stationary, that no
small ``virtual'' perturbation $v$ of the field changes the energy to first
order. We will make this precise in \cref{sec:virtual-work}; for now, hold onto
the slogan.

\begin{keyidea}
A weak form replaces ``the equation holds at every point'' with ``the equation,
averaged against every test function, holds.'' Integration by parts then trades
one derivative on the unknown for one derivative on the test function. Two gains
follow at once: the solution may be \emph{less smooth} (one derivative instead
of two), and the equation takes the form of a \emph{variational} (energy)
principle---the form in which the physics naturally lives.
\end{keyidea}

The rest of this chapter executes this move three times, with increasing
generality, and then extracts the single template they share.

\section{The one-dimensional archetype}
\label{sec:1d}

We begin where every student of the subject begins: the one-dimensional
Poisson equation. It is small enough to do completely by hand, and every
structural feature of the general theory is already present in it.

Let the domain be the interval $\dom = (0,1)$, and consider the boundary-value
problem
\begin{equation}
  -u''(x) = f(x) \quad\text{for } x\in(0,1), \qquad u(0)=u(1)=0.
  \label{eq:1d-strong}
\end{equation}
Here $u$ is the unknown (think of it as a temperature along an insulated rod, or
the deflection of a stretched string under a transverse load $f$), and the
boundary conditions pin both ends to zero. This is the strong form: it asks for
a function $u$, twice differentiable, whose second derivative equals $-f$
everywhere.

\subsection{Multiply, integrate, integrate by parts}

Choose a \emph{test function} $v$. For now we require only two things of it: it
should be once-differentiable (so that $v'$ makes sense), and---this is the
crucial choice---it should \emph{vanish at the boundary}, $v(0)=v(1)=0$. We will
see in a moment exactly why. Multiply \cref{eq:1d-strong} by $v$ and integrate
over $(0,1)$:
\begin{equation}
  -\int_0^1 u''(x)\,v(x)\dx = \int_0^1 f(x)\,v(x)\dx.
  \label{eq:1d-weighted}
\end{equation}
Now integrate the left-hand side by parts. Recall the one-dimensional formula
$\int_0^1 p'\,q = \bigl[p\,q\bigr]_0^1 - \int_0^1 p\,q'$. Apply it with $p=u'$
(so $p'=u''$) and $q=v$:
\begin{equation}
  -\int_0^1 u''\,v\dx
  = -\Bigl[\,u'(x)\,v(x)\,\Bigr]_0^1 + \int_0^1 u'(x)\,v'(x)\dx.
  \label{eq:1d-ibp}
\end{equation}
Look at the boundary term $\bigl[u'\,v\bigr]_0^1 = u'(1)v(1) - u'(0)v(0)$. This
is exactly where our requirement $v(0)=v(1)=0$ earns its keep: \emph{both}
$v(1)$ and $v(0)$ are zero, so the entire boundary term vanishes. We chose the
test functions to die at the boundary precisely so that this term would
disappear. What survives is clean:
\begin{equation}
  \int_0^1 u'(x)\,v'(x)\dx = \int_0^1 f(x)\,v(x)\dx
  \qquad\text{for all admissible } v.
  \label{eq:1d-weak}
\end{equation}
This is the \emph{weak form} of \cref{eq:1d-strong}. Read it carefully and
notice what has changed. On the left, $u$ now appears only through its
\emph{first} derivative $u'$, never $u''$. The unknown no longer needs to be
twice differentiable; it needs one derivative, and that derivative may be
merely square-integrable. The price was symmetry: where the strong form treated
$u$ specially, the weak form treats $u$ and $v$ even-handedly---each contributes
exactly one derivative.

\begin{figure}[t]
\centering
\begin{tikzpicture}[scale=1.0]
  % axis
  \draw[->,simGray] (-0.3,0) -- (5.3,0) node[right,font=\scriptsize]{$x$};
  \draw[->,simGray] (0,-0.3) -- (0,2.4);
  \node[font=\scriptsize,simGray,below] at (0,0) {$0$};
  \node[font=\scriptsize,simGray,below] at (5,0) {$1$};
  \draw[simGray,densely dotted] (5,0) -- (5,2.2);
  % the solution u (a smooth bump)
  \draw[simBlue,very thick,domain=0:5,samples=80]
    plot (\x,{1.9*sin(180*\x/5)});
  \node[simBlue,font=\small] at (4.0,1.95) {$u(x)$};
  % a hat test function v
  \draw[simRust,thick] (0,0) -- (2.4,1.25) -- (5,0);
  \node[simRust,font=\small] at (1.9,1.45) {$v(x)$};
  % endpoints clamped
  \fill[simInk] (0,0) circle (1.6pt);
  \fill[simInk] (5,0) circle (1.6pt);
  \node[font=\scriptsize,simRust] at (3.6,0.55) {$v(0)=v(1)=0$};
\end{tikzpicture}
\caption[A test function against the unknown]{The unknown $u$ (blue) and one
admissible test function $v$ (rust), here a piecewise-linear ``hat.'' The test
function is required to vanish at both ends, $v(0)=v(1)=0$; this is what kills
the boundary term in the integration by parts. A hat like this is the prototype
of the finite-element basis functions of \cref{ch:functionspaces}: the weak form
only ever pairs $u'$ against $v'$, so $v$ needs just one (possibly jumpy)
derivative.}
\label{fig:hat}
\end{figure}

\subsection{Naming the parts: bilinear form and linear functional}

The weak form \cref{eq:1d-weak} has a structure worth naming, because the names
will persist unchanged through every generalization to come. Define
\begin{align}
  a(u,v) &\;=\; \int_0^1 u'(x)\,v'(x)\dx, \label{eq:1d-bilinear}\\
  \ell(v) &\;=\; \int_0^1 f(x)\,v(x)\dx. \label{eq:1d-linear}
\end{align}
The first, $a(u,v)$, is a \emph{bilinear form}: it takes two functions and
returns a number, and it is linear in each argument separately. Linearity in
the first slot means $a(\alpha u_1 + \beta u_2, v) = \alpha\,a(u_1,v) +
\beta\,a(u_2,v)$, and similarly in the second---both are immediate from the
linearity of the integral and the derivative. The second, $\ell(v)$, is a
\emph{linear functional}: it takes one function and returns a number, linearly.
With these names, the entire weak problem collapses to one line:
\begin{equation}
  \boxed{\;\text{find } u \text{ with } u(0)=u(1)=0 \text{ such that } a(u,v)=\ell(v) \text{ for all admissible } v.\;}
  \label{eq:1d-abstract}
\end{equation}
This abstract shape---$a(u,v) = \ell(v)$, a symmetric bilinear form equated to a
linear functional---is the form the finite element method discretizes. When in
\cref{ch:galerkin} we replace the infinite-dimensional space of admissible
functions with a finite-dimensional one spanned by hat functions like the one in
\cref{fig:hat}, the bilinear form $a$ becomes a matrix, the functional $\ell$
becomes a vector, and \cref{eq:1d-abstract} becomes a linear system $\mat{K}\vx
= \vb$. Everything downstream is bookkeeping on $a$ and $\ell$.

\begin{notationbox}
We will write the bilinear form as an inner product whenever it is convenient:
$a(u,v) = \inner{u'}{v'}$, where $\inner{\cdot}{\cdot}$ is the $\Ltwo(\dom)$
inner product $\inner{p}{q} = \int_\dom p\,q$. The two notations
$\int_0^1 u'v'\dx$ and $\inner{u'}{v'}$ mean the same thing; the inner-product
notation makes the structure ``differential operator applied to each argument,
paired in $\Ltwo$'' visible at a glance, and it is the notation in which the
unifying family of \cref{sec:family} is cleanest.
\end{notationbox}

\subsection{A complete worked example}

Let us carry the 1-D problem all the way to a number, so that nothing is
abstract.

\begin{workedexample}{The 1-D weak form with a concrete load}{1dfull}
Solve, in weak form, the problem
\[
  -u''(x) = \pi^2\sin(\pi x) \quad\text{on } (0,1), \qquad u(0)=u(1)=0,
\]
and verify the weak statement on a specific test function.

\medskip\noindent\textbf{Strong solution (for checking).}
The strong problem has the exact solution $u(x)=\sin(\pi x)$: indeed
$u''=-\pi^2\sin(\pi x)$, so $-u'' = \pi^2\sin(\pi x)=f$, and $u(0)=u(1)=0$. We
keep this in our pocket to check the weak form against.

\medskip\noindent\textbf{The bilinear form and functional.}
With $f(x)=\pi^2\sin(\pi x)$ the weak problem is: find $u$ with $u(0)=u(1)=0$ such
that, for every admissible $v$,
\[
  a(u,v)=\int_0^1 u'(x)\,v'(x)\dx
  \;=\;
  \ell(v)=\int_0^1 \pi^2\sin(\pi x)\,v(x)\dx.
\]

\medskip\noindent\textbf{Test it on $v(x)=\sin(\pi x)$.}
This $v$ is admissible: it is smooth and vanishes at both ends. With the candidate
solution $u(x)=\sin(\pi x)$ we have $u'(x)=\pi\cos(\pi x)$ and
$v'(x)=\pi\cos(\pi x)$, so the left side is
\[
  a(u,v)=\int_0^1 \pi^2\cos^2(\pi x)\dx
  = \pi^2\cdot\tfrac12 = \tfrac{\pi^2}{2},
\]
using $\int_0^1\cos^2(\pi x)\dx=\tfrac12$. The right side is
\[
  \ell(v)=\int_0^1 \pi^2\sin^2(\pi x)\dx
  = \pi^2\cdot\tfrac12 = \tfrac{\pi^2}{2},
\]
using $\int_0^1\sin^2(\pi x)\dx=\tfrac12$. The two agree: $a(u,v)=\ell(v)=\pi^2/2$.
The weak equation is satisfied on this test function, as it must be for the true
solution.

\medskip\noindent\textbf{Where the boundary term went.}
The integration by parts produced the term $-[u'v]_0^1 = -\bigl(u'(1)v(1) -
u'(0)v(0)\bigr)$. Here $u'(x)=\pi\cos(\pi x)$, so $u'(1)=-\pi$ and $u'(0)=\pi$ are
both \emph{nonzero}; the boundary term vanished entirely because $v(1)=v(0)=0$,
not because the derivative of $u$ was small. This is the general mechanism: the
homogeneous Dirichlet condition on the \emph{test} space, not any property of the
solution, is what discards the boundary contribution.
\end{workedexample}

\subsection{What changes at a Neumann end}
\label{sec:1d-neumann}

The Dirichlet condition $u(0)=u(1)=0$ in \cref{eq:1d-strong} entered the weak
form in two distinct ways, and it is worth separating them, because the
distinction is the seed of the whole essential-versus-natural story of
\cref{sec:essential-natural}. First, the condition on $u$ itself restricts the
\emph{trial space}: we look for the solution among functions that already vanish
at the ends. Second, the matching condition on $v$ restricts the \emph{test
space}, and that is what annihilated the boundary term. The boundary condition on
$u$ is enforced by \emph{constraining the space}; the boundary term, meanwhile,
is a separate object that we made disappear by a choice on $v$.

Now change one end. Suppose instead of $u(1)=0$ we are told the \emph{derivative}
there: $u'(1) = g$ for some given flux $g$ (a Neumann, or natural, condition---
physically, a prescribed heat flux or surface charge), while keeping $u(0)=0$.
The strong problem is now
\[
  -u'' = f \text{ on } (0,1), \qquad u(0)=0, \quad u'(1)=g.
\]
Repeat the derivation. We multiply by $v$ and integrate by parts exactly as
before, arriving again at \cref{eq:1d-ibp}. But now we can no longer ask $v(1)=0$:
the value of $u$ at $x=1$ is \emph{not} prescribed, so the trial space does not
constrain it, and correspondingly we must leave the test functions free there.
We retain $v(0)=0$ (the left end is still Dirichlet). The boundary term is
therefore
\[
  -\Bigl[u'v\Bigr]_0^1 = -\bigl(u'(1)v(1) - u'(0)v(0)\bigr)
  = -u'(1)\,v(1) = -g\,v(1),
\]
because $v(0)=0$ kills the left contribution while $u'(1)=g$ is \emph{known}. The
weak form becomes
\begin{equation}
  \underbrace{\int_0^1 u'v'\dx}_{a(u,v)}
  \;=\;
  \underbrace{\int_0^1 f\,v\dx + g\,v(1)}_{\ell(v)}.
  \label{eq:1d-neumann-weak}
\end{equation}
Two things happened, and both are general. The Dirichlet end ($x=0$) was
enforced by constraining the space and produced nothing in the equation---it is
an \emph{essential} condition. The Neumann end ($x=1$) was \emph{not} imposed on
the space at all; instead the prescribed flux $g$ rode in through the surviving
boundary term and landed in the functional $\ell$ as a source---it is a
\emph{natural} condition. A homogeneous Neumann condition ($g=0$) would land
nothing and need no enforcement: it is the condition you get ``for free'' by
doing nothing, which is why the weak form is sometimes said to satisfy it
naturally. Keep these two roles in view; \cref{sec:essential-natural} is nothing
but their three-dimensional restatement.

\section{The three-dimensional scalar case: electrostatics}
\label{sec:3d-scalar}

We now lift the 1-D archetype to three dimensions. The structure will be
identical; only the integration-by-parts identity changes, from the
one-dimensional formula to the divergence theorem. This is the case that
governs the electrostatic driver---the capacitance analysis of
\cref{ch:targets}---so we work it concretely.

The strong form is \cref{eq:strong-es},
\[
  -\Div(\varepsilon\,\Grad u) = f \quad\text{in } \dom,
\]
with $\dom\subset\Real^3$ a bounded domain, $\varepsilon$ the (possibly
position-dependent, possibly matrix-valued) permittivity, $u$ the electrostatic
potential, and $f$ a charge density. Following the 1-D recipe to the letter:
choose a scalar test function $v$, multiply, and integrate over $\dom$:
\begin{equation}
  -\int_\dom \Div(\varepsilon\Grad u)\,v\dV = \int_\dom f\,v\dV.
  \label{eq:3d-weighted}
\end{equation}

\subsection{Green's first identity}

The three-dimensional analogue of integration by parts is \emph{Green's first
identity}, and it is itself a one-line consequence of the divergence theorem. We
derive it so that nothing is taken on faith. Start from the product rule for the
divergence of a scalar times a vector field. For a scalar field $v$ and a vector
field $\vF$,
\begin{equation}
  \Div(v\,\vF) = \Grad v\cdot\vF + v\,\Div\vF.
  \label{eq:prodrule}
\end{equation}
Integrate this over $\dom$ and apply the divergence theorem
$\int_\dom\Div(\,\cdot\,)\dV=\oint_\bdry(\,\cdot\,)\cdot\vn\,\mathrm{d}S$ to the
left side, where $\vn$ is the outward unit normal on the boundary $\bdry$:
\[
  \oint_\bdry v\,\vF\cdot\vn\,\mathrm{d}S
  = \int_\dom \Grad v\cdot\vF\dV + \int_\dom v\,\Div\vF\dV.
\]
Rearranging,
\begin{equation}
  \int_\dom v\,\Div\vF\dV
  = -\int_\dom \Grad v\cdot\vF\dV
  + \oint_\bdry v\,(\vF\cdot\vn)\,\mathrm{d}S.
  \label{eq:green1}
\end{equation}
This is Green's first identity. Compare it term-for-term with the 1-D formula
\cref{eq:1d-ibp}: $\int v\,\Div\vF$ plays the role of $\int u''v$, the
volume term $\int\Grad v\cdot\vF$ is the new $\int u'v'$, and the surface
integral $\oint_\bdry$ is the boundary term $[u'v]_0^1$. The divergence theorem
\emph{is} integration by parts in three dimensions; the only genuine change is
that the ``boundary'' of an interval (two points) has become the boundary of a
solid (a surface).

\subsection{Applying it to the electrostatic problem}

Set $\vF = \varepsilon\Grad u$ in Green's identity \cref{eq:green1}. Then
$\Div\vF = \Div(\varepsilon\Grad u)$, and the identity reads
\[
  \int_\dom v\,\Div(\varepsilon\Grad u)\dV
  = -\int_\dom \Grad v\cdot(\varepsilon\Grad u)\dV
  + \oint_\bdry v\,(\varepsilon\Grad u)\cdot\vn\,\mathrm{d}S.
\]
Substitute this into the left side of the weighted residual
\cref{eq:3d-weighted} (which carries an overall minus sign):
\begin{equation}
  \int_\dom \varepsilon\,\Grad u\cdot\Grad v\dV
  \;-\;
  \oint_\bdry v\,(\varepsilon\Grad u)\cdot\vn\,\mathrm{d}S
  \;=\;
  \int_\dom f\,v\dV.
  \label{eq:3d-with-boundary}
\end{equation}
There is the volume bilinear form we were after, $\int_\dom\varepsilon\Grad
u\cdot\Grad v$, sitting on the left exactly as $\int u'v'$ did in one dimension.
And there is the boundary integral, the three-dimensional descendant of the
$[u'v]_0^1$ term, carrying the normal flux $(\varepsilon\Grad u)\cdot\vn$ across
$\bdry$. What we do with that surface integral---make it vanish, or keep it as a
source---is precisely the essential/natural dichotomy, and it is important
enough to get its own section.

\subsection{Splitting the boundary: essential and natural}
\label{sec:essential-natural}

Partition the boundary into two disjoint pieces,
\[
  \bdry = \Gamma_D \cup \Gamma_N, \qquad \Gamma_D\cap\Gamma_N=\emptyset,
\]
where on $\Gamma_D$ (the \emph{Dirichlet} part) the potential is prescribed,
$u = u_0$, and on $\Gamma_N$ (the \emph{Neumann} part) the normal flux is
prescribed, $(\varepsilon\Grad u)\cdot\vn = g$. Physically, $\Gamma_D$ is a
conductor held at a fixed voltage and $\Gamma_N$ is a surface with a given
surface charge (with $g=0$ the natural ``no-flux'' wall).

\begin{figure}[t]
\centering
\begin{tikzpicture}[scale=1.0]
  % a blobby domain
  \fill[simBgBlue]
    (0,0) .. controls (0.4,1.6) and (2.2,2.0) .. (3.6,1.4)
          .. controls (4.8,0.95) and (4.6,-0.6) .. (3.2,-1.0)
          .. controls (1.8,-1.4) and (-0.4,-1.0) .. (0,0) -- cycle;
  % Dirichlet portion (thick blue arc, lower-left)
  \draw[simBlue,line width=2.2pt]
    (0,0) .. controls (1.8,-1.4) and (3.2,-1.0) .. (3.2,-1.0);
  % Neumann portion (thick green arc, upper)
  \draw[simGreen,line width=2.2pt]
    (0,0) .. controls (0.4,1.6) and (2.2,2.0) .. (3.6,1.4)
          .. controls (4.8,0.95) and (4.6,-0.6) .. (3.2,-1.0);
  \node[simInk,font=\large] at (2.1,0.3) {$\dom$};
  \node[simBlue,font=\small] at (1.4,-1.5) {$\Gamma_D:\ u=u_0$};
  \node[simGreen,font=\small] at (4.7,1.55) {$\Gamma_N:\ (\varepsilon\Grad u)\cdot\vn=g$};
  % outward normal on Neumann part
  \draw[-{Stealth[length=2mm]},simGreen,thick] (3.55,1.42) -- (4.25,1.85)
     node[right,font=\scriptsize]{$\vn$};
  % test function vanishes on Dirichlet
  \node[simBlue,font=\scriptsize\itshape] at (1.5,-0.75) {$v=0$ here};
\end{tikzpicture}
\caption[The boundary splits into Dirichlet and Neumann parts]{Green's first
identity produces a boundary integral over all of $\bdry$. Splitting
$\bdry=\Gamma_D\cup\Gamma_N$ disposes of it in two ways. On the Dirichlet part
$\Gamma_D$ (blue) the potential is prescribed and we constrain the test space to
vanish, $v=0$, so the integral there drops---the condition is \emph{essential}.
On the Neumann part $\Gamma_N$ (green) the normal flux $(\varepsilon\Grad
u)\cdot\vn=g$ is prescribed; the integral survives and carries $g$ into the
right-hand side as a surface source---the condition is \emph{natural}.}
\label{fig:greenbdry}
\end{figure}

Now dispose of the boundary integral $\oint_\bdry v(\varepsilon\Grad
u)\cdot\vn\,\mathrm{d}S$ piece by piece, exactly as in the 1-D Neumann
discussion of \cref{sec:1d-neumann}.

\emph{On $\Gamma_D$}, where $u$ is prescribed, we have no information about the
flux $(\varepsilon\Grad u)\cdot\vn$---it is part of the unknown response. We
remove the difficulty the same way we did in one dimension: we demand that the
test functions vanish on $\Gamma_D$, $v|_{\Gamma_D}=0$. Then the integral over
$\Gamma_D$ is zero, regardless of the unknown flux. The Dirichlet datum $u=u_0$
is enforced separately, by constraining the trial space to functions that equal
$u_0$ on $\Gamma_D$. This is the \emph{essential} boundary condition: it lives in
the \emph{space}, not in the equation, and it leaves no trace in $a$ or $\ell$.

\emph{On $\Gamma_N$}, where the flux is prescribed, the integrand
$(\varepsilon\Grad u)\cdot\vn = g$ is \emph{known}. We keep the test functions
free there, substitute $g$, and move the resulting surface integral to the
right-hand side. It becomes a source term in the functional. This is the
\emph{natural} boundary condition: it lives in the \emph{equation}, contributing
to $\ell$.

Putting it together, the electrostatic weak form is
\begin{equation}
  \boxed{\;
  \underbrace{\int_\dom \varepsilon\,\Grad u\cdot\Grad v\dV}_{a(u,v)}
  \;=\;
  \underbrace{\int_\dom f\,v\dV \;+\; \oint_{\Gamma_N} g\,v\,\mathrm{d}S}_{\ell(v)}
  \;}
  \label{eq:3d-weak}
\end{equation}
for all test functions $v$ with $v|_{\Gamma_D}=0$, sought among trial functions
$u$ with $u|_{\Gamma_D}=u_0$. The homogeneous-Neumann case $g=0$---the most
common wall condition in a capacitance problem, where no charge sits on the
insulating outer boundary---drops the surface integral entirely and is enforced
by doing nothing at all. That is the sense in which Neumann conditions are
``natural.''

\begin{pitfall}
It is tempting to think the essential (Dirichlet) condition is the ``easy'' one
because it looks simpler in the strong form. In the weak/discrete setting it is
the \emph{harder} one to implement: it must be built into the function space, by
eliminating or constraining the degrees of freedom that sit on $\Gamma_D$
(\cref{ch:bc}). The natural (Neumann) condition, which looks like the
complicated one, requires almost nothing---a homogeneous Neumann wall is handled
by simply not assembling a boundary term. ``Essential'' means ``you must enforce
it yourself''; ``natural'' means ``the weak form enforces it for you.''
\end{pitfall}

\begin{workedexample}{The electrostatic weak form, term by term}{esweak}
Take a parallel-plate capacitor: $\dom$ the slab between two conducting plates,
the top plate held at $u=V$ and the bottom at $u=0$ (both Dirichlet, so both part
of $\Gamma_D$), the side walls insulating (homogeneous Neumann, $g=0$, so part of
$\Gamma_N$), no free charge in the dielectric ($f=0$), and a uniform scalar
permittivity $\varepsilon$.

\medskip\noindent\textbf{Identify the pieces of the template.}
Reading \cref{eq:3d-weak} against $a(u,v)=(Q\,\diffop u,\diffop v)$:
\begin{itemize}\setlength\itemsep{1pt}
  \item the coefficient is $Q=\varepsilon$, the permittivity;
  \item the differential operator is $\diffop=\Grad$, the gradient;
  \item the bilinear form is $a(u,v)=\inner{\varepsilon\Grad u}{\Grad v}_{\Ltwo(\dom)}=\int_\dom\varepsilon\,\Grad u\cdot\Grad v\dV$.
\end{itemize}

\medskip\noindent\textbf{Assemble the functional.}
With $f=0$ in the volume and $g=0$ on the insulating walls, the functional is
\[
  \ell(v) = \int_\dom 0\cdot v\dV + \oint_{\Gamma_N} 0\cdot v\,\mathrm{d}S = 0.
\]
The right-hand side is \emph{empty}. The problem is not, however, trivial: the
nonzero plate voltage $u=V$ is an \emph{inhomogeneous essential} condition, and
it enters not through $\ell$ but through the trial space (the solution must equal
$V$ on the top plate). In the discrete setting this inhomogeneous Dirichlet datum
is ``lifted'' into the right-hand side by the elimination procedure of
\cref{ch:bc}---which is exactly why an apparently source-free problem
($f=g=0$) still produces a nonzero load vector. The physics confirms it: a
charged capacitor with no free charge in the dielectric still stores energy,
driven entirely by the imposed plate voltages.

\medskip\noindent\textbf{The product.}
Solving the weak problem yields the potential $u$; the capacitance follows from
the stored energy $U=\tfrac12\int_\dom\varepsilon\abs{\Grad u}^2\dV =
\tfrac12\,a(u,u)$ via $C = 2U/V^2$. The bilinear form $a$ evaluated on the
solution against \emph{itself} is the energy---a connection we exploit in
\cref{sec:virtual-work} and again in the capacitance reduction of
\cref{ch:electrostatics}.
\end{workedexample}

\section{The vector case: curl--curl}
\label{sec:curlcurl}

The electrostatic problem lives on a \emph{scalar} field, the potential $u$, and
its natural derivative is the gradient. The magnetostatic and eigenmode problems
of \cref{ch:reductions-pde} live on \emph{vector} fields---the magnetic vector
potential $\vA$, the electric field $\vE$---and their natural derivative is the
\emph{curl}. The integration-by-parts move is structurally identical, but the
identity it rests on is the curl analogue of Green's formula, and the function
space is different. We sketch the derivation and read off the same template.

The strong form of the magnetostatic problem is the curl--curl equation
\begin{equation}
  \Curl(\nu\,\Curl\vu) = \vJ \quad\text{in }\dom,
  \label{eq:curlcurl-strong}
\end{equation}
where $\vu$ is the unknown vector field (the vector potential $\vA$), $\nu=1/\mu$
is the reluctivity (the magnetic coefficient, the analogue of $\varepsilon$), and
$\vJ$ is the source current density. Multiply by a vector test function $\vv$ and
integrate---now the multiplication is a dot product, since both factors are
vectors:
\begin{equation}
  \int_\dom \Curl(\nu\Curl\vu)\cdot\vv\dV = \int_\dom \vJ\cdot\vv\dV.
  \label{eq:curlcurl-weighted}
\end{equation}

\subsection{The curl integration-by-parts identity}

The identity we need plays the exact role Green's first identity played for the
gradient. For two vector fields $\vp$ and $\vv$,
\begin{equation}
  \int_\dom (\Curl\vp)\cdot\vv\dV
  = \int_\dom \vp\cdot(\Curl\vv)\dV
  + \oint_\bdry (\vp\times\vn)\cdot\vv\,\mathrm{d}S.
  \label{eq:curlibp}
\end{equation}
Like Green's identity, this follows from a vector product rule---$\Div(\vp\times
\vv) = (\Curl\vp)\cdot\vv - \vp\cdot(\Curl\vv)$---integrated over $\dom$ with the
divergence theorem applied to the left side; the surface term
$\oint_\bdry(\vp\times\vv)\cdot\vn\,\mathrm{d}S$ is rewritten as
$\oint_\bdry(\vp\times\vn)\cdot\vv\,\mathrm{d}S$ by the scalar triple-product
identity. The shape is the same as before: a derivative on $\vp$ moves to a
derivative on $\vv$, at the cost of a boundary integral. Apply it with
$\vp = \nu\Curl\vu$:
\[
  \int_\dom \Curl(\nu\Curl\vu)\cdot\vv\dV
  = \int_\dom (\nu\Curl\vu)\cdot(\Curl\vv)\dV
  + \oint_\bdry \bigl((\nu\Curl\vu)\times\vn\bigr)\cdot\vv\,\mathrm{d}S.
\]

\begin{figure}[t]
\centering
\begin{tikzpicture}[
    row/.style={font=\small},
    lbl/.style={font=\footnotesize\bfseries, text=simInk},
  ]
  % column headers
  \node[lbl] at (0,2.0) {dimension / operator};
  \node[lbl] at (3.7,2.0) {move a derivative\dots};
  \node[lbl] at (8.4,2.0) {\dots leaving a boundary term};
  \draw[simGray,semithick] (-2.0,1.6) -- (10.6,1.6);
  % 1-D row
  \node[row,simBlue,anchor=west] at (-2.0,1.0)
    {1-D: \ $\displaystyle\int u''\,v$};
  \node[row,anchor=west] at (2.3,1.0)
    {$\displaystyle =-\int u'\,v'$};
  \node[row,simRust,anchor=west] at (6.6,1.0)
    {$+\,\bigl[u'\,v\bigr]_0^1$};
  % 3-D gradient row
  \node[row,simBlue,anchor=west] at (-2.0,0.1)
    {$\Grad$: \ $\displaystyle\int v\,\Div\vF$};
  \node[row,anchor=west] at (2.3,0.1)
    {$\displaystyle =-\int \Grad v\cdot\vF$};
  \node[row,simRust,anchor=west] at (6.6,0.1)
    {$+\displaystyle\oint_\bdry v\,(\vF\cdot\vn)$};
  % curl row
  \node[row,simBlue,anchor=west] at (-2.0,-0.8)
    {$\Curl$: \ $\displaystyle\int(\Curl\vp)\cdot\vv$};
  \node[row,anchor=west] at (2.3,-0.8)
    {$\displaystyle =\int \vp\cdot(\Curl\vv)$};
  \node[row,simRust,anchor=west] at (6.6,-0.8)
    {$+\displaystyle\oint_\bdry(\vp\times\vn)\cdot\vv$};
  \draw[simGray,semithick] (-2.0,-1.3) -- (10.6,-1.3);
  \node[font=\scriptsize\itshape,simGray,anchor=west] at (-2.0,-1.75)
    {one move, three settings: a derivative crosses from one factor to the other, and a boundary term is born.};
\end{tikzpicture}
\caption[The three integration-by-parts identities, aligned]{The same move in
three settings. One-dimensional integration by parts, Green's first identity (for
the gradient), and the curl integration-by-parts identity all have the identical
shape: a derivative is transferred from one factor to the other (middle column),
at the cost of a boundary integral (rust). The 1-D boundary ``surface'' is two
endpoints; in three dimensions it is a genuine surface $\bdry$. Choosing test
functions that vanish (gradient case) or whose tangential part vanishes (curl
case) on the Dirichlet boundary kills that term---the one mechanism that recurs
throughout the chapter.}
\label{fig:threeibp}
\end{figure}

Substituting into the weighted residual \cref{eq:curlcurl-weighted} and, as
before, choosing test functions whose \emph{tangential} component vanishes on the
Dirichlet boundary (so the surface term drops there) gives the weak form
\begin{equation}
  \boxed{\;
  \underbrace{\int_\dom \nu\,(\Curl\vu)\cdot(\Curl\vv)\dV}_{a(\vu,\vv)}
  \;=\;
  \underbrace{\int_\dom \vJ\cdot\vv\dV \;+\;(\text{natural-BC surface term})}_{\ell(\vv)}
  \;}
  \label{eq:curlcurl-weak}
\end{equation}
for all admissible $\vv$. The natural boundary condition here prescribes the
\emph{tangential} field $\nu\Curl\vu\times\vn$ (physically, a tangential
magnetic field $\vH$); the essential condition prescribes the tangential
component of $\vu$ itself. The pattern is the gradient story with ``value on the
boundary'' replaced by ``tangential component on the boundary''---a change forced
by the geometry of the curl, treated fully in \cref{ch:functionspaces} and
\cref{ch:bc}.

\subsection{The mass term}
\label{sec:mass}

Eigenmode and driven problems pair the curl--curl stiffness term with a second,
\emph{lower-order} term that carries no derivative at all. In the eigenmode
problem $\Curl(\nu\Curl\vE)=\omega^2\varepsilon\vE$, the right side
$\omega^2\varepsilon\vE$ contributes, after multiplying by $\vv$ and integrating,
the term
\begin{equation}
  m(\vu,\vv) = \int_\dom \varepsilon\,\vu\cdot\vv\dV,
  \label{eq:mass}
\end{equation}
with \emph{no} integration by parts---there is no derivative to move. This is the
\emph{mass} term (a name inherited from the analogous term in structural
dynamics, where the coefficient is literally a mass density). It is the
``identity'' member of the family: the differential operator applied to $\vu$ and
$\vv$ is the identity, $\diffop\vu=\vu$, and the pairing is just the
$Q$-weighted $\Ltwo$ inner product. The eigenmode weak form is then the
\emph{pencil} $a(\vu,\vv)=\omega^2\,m(\vu,\vv)$, which discretizes to the
generalized eigenvalue problem $\mat{K}\vx=\lambda\mat{M}\vx$ of
\cref{ch:targets}---the stiffness $\mat{K}$ from the curl--curl term, the mass
$\mat{M}$ from \cref{eq:mass}.

That the curl--curl term and the mass term are \emph{the same kind of object}---a
coefficient-weighted $\Ltwo$ pairing of a differential operator applied to both
arguments, differing only in \emph{which} operator---is the observation the next
section turns into the chapter's payoff.

\section{The unifying bilinear-form family}
\label{sec:family}

Lay the three bilinear forms we have derived side by side, written in the
inner-product notation of the notation box:
\begin{align*}
  \text{electrostatic (gradient):}\quad & a(u,v) = \inner{\varepsilon\,\Grad u}{\Grad v}_{\Ltwo(\dom)},\\
  \text{magnetostatic (curl):}\quad & a(\vu,\vv) = \inner{\nu\,\Curl\vu}{\Curl\vv}_{\Ltwo(\dom)},\\
  \text{mass (identity):}\quad & m(\vu,\vv) = \inner{\varepsilon\,\vu}{\vv}_{\Ltwo(\dom)}.
\end{align*}
The three are visibly instances of one shape. In each, a differential operator
is applied to both the trial and the test function; the two results are paired in
the $\Ltwo(\dom)$ inner product; and a material coefficient $Q$ weights the
pairing. Only \emph{which differential operator} changes: $\Grad$, then $\Curl$,
then the identity. Naming the differential operator $\diffop$ and the coefficient
$Q$, all three collapse into a single template.

\begin{keyidea}
\textbf{The bilinear-form family.}
Every domain bilinear form the simulator assembles has the form
\begin{equation}
  a(u,v) \;=\; \bigl(Q\,\diffop u,\ \diffop v\bigr)_{\Ltwo(\dom)}
            \;=\; \int_\dom (Q\,\diffop u)\cdot(\diffop v)\dV,
  \label{eq:family}
\end{equation}
where $\diffop\in\{\,\text{value},\ \Grad,\ \Curl,\ \Div\,\}$ is a differential
operator applied to both trial and test functions, and $Q$ is a material
coefficient (scalar or matrix-valued). A weak-form problem is built by choosing,
for each term, a pair $(Q,\diffop)$. The four choices of $\diffop$ are:
\[
\begin{array}{llll}
  \diffop=\text{value} & a(u,v)=(Q\,u,\,v) & \text{mass / }\Ltwo\text{ pairing} & \text{any space}\\[2pt]
  \diffop=\Grad        & a(u,v)=(Q\,\Grad u,\,\Grad v) & \text{electrostatic (diffusion)} & \Hone\\[2pt]
  \diffop=\Curl        & a(u,v)=(Q\,\Curl u,\,\Curl v) & \text{magnetostatic / eigenmode} & \Hcurl\\[2pt]
  \diffop=\Div         & a(u,v)=(Q\,\Div u,\,\Div v) & \text{div--div} & \Hdiv
\end{array}
\]
\end{keyidea}

This template is the load-bearing abstraction of the entire field-theory part of
the book, and it is worth pausing on what it buys us. A weak-form term is now a
\emph{piece of data}: a pair $(\text{coefficient } Q,\ \text{differential
operator } \diffop)$. To specify the physics of a problem is to write down a
short list of such pairs. The electrostatic stiffness is the single term
$(\varepsilon,\Grad)$; the magnetostatic stiffness is $(\nu,\Curl)$; the
eigenmode problem is two terms, $(\nu,\Curl)$ for the stiffness and
$(\varepsilon,\text{value})$ for the mass. The machinery that turns such a pair
into an actual matrix---element quadrature, basis evaluation, assembly---is
\emph{completely indifferent} to which pair it was handed; it is the same fold
over terms in every case, the subject of \cref{ch:galerkin} and
\cref{ch:assembly}.

\begin{figure}[t]
\centering
\setlength{\tabcolsep}{0pt}
\renewcommand{\arraystretch}{1.0}
\begin{tikzpicture}[
    cell/.style={draw=simGray, semithick, minimum width=33mm, minimum height=20mm,
      align=center, font=\footnotesize, inner sep=2pt},
    hdr/.style={font=\small\bfseries, text=simInk},
  ]
  % four cells in a row
  \node[cell, fill=simBgGray] (val) at (0,0)
    {$\diffop=\text{value}$\\[3pt]$(Q\,u,\;v)$\\[3pt]\textit{mass}\\[1pt]\scriptsize any space};
  \node[cell, fill=simBgBlue, right=2mm of val] (grad)
    {$\diffop=\Grad$\\[3pt]$(Q\,\Grad u,\;\Grad v)$\\[3pt]\textit{electrostatic}\\[1pt]\scriptsize $\Hone$};
  \node[cell, fill=simBgRust, right=2mm of grad] (curl)
    {$\diffop=\Curl$\\[3pt]$(Q\,\Curl u,\;\Curl v)$\\[3pt]\textit{magnetostatic}\\[1pt]\scriptsize $\Hcurl$};
  \node[cell, fill=simBgGreen, right=2mm of curl] (div)
    {$\diffop=\Div$\\[3pt]$(Q\,\Div u,\;\Div v)$\\[3pt]\textit{div--div}\\[1pt]\scriptsize $\Hdiv$};
  % unifying brace above
  \draw[decorate,decoration={brace,amplitude=5pt},simInk]
    ($(val.north west)+(0,3mm)$) -- ($(div.north east)+(0,3mm)$)
    node[midway,above=5pt,font=\small\bfseries]{$a(u,v)=(Q\,\diffop u,\;\diffop v)_{\Ltwo(\dom)}$};
  % de Rham hint below
  \node[font=\scriptsize\itshape,simGray] at ($(val.south)!0.5!(div.south)+(0,-6mm)$)
    {one template, four differential operators --- the de Rham ladder of \cref{ch:derham}};
\end{tikzpicture}
\caption[The four-operator bilinear-form family]{The bilinear-form family as a
labeled grid. A single template $a(u,v)=(Q\,\diffop u,\diffop v)_{\Ltwo(\dom)}$
(brace) specializes to four terms by the choice of differential operator
$\diffop$. The value (mass) term carries no derivative; gradient, curl, and
divergence each carry one, and each is well-formed on a different
finite-element space---$\Hone$, $\Hcurl$, $\Hdiv$---the slots of the de Rham
complex of \cref{ch:derham}. The coefficient $Q$ rides every cell uniformly.}
\label{fig:family}
\end{figure}

\begin{remark}[How Palace records a term]
This is not an abstraction imposed from outside; it is the structure Palace's own
code already has. The simulator builds an operator by pushing
\emph{integrator} objects onto a list, each integrator carrying a coefficient and
a differential-operator type: the diffusion integrator for $\diffop=\Grad$, the
curl--curl integrator for $\diffop=\Curl$, the mass integrator for the value
term. Every integrator inherits the same coefficient slot from a shared base
class, which is exactly the statement that $Q$ is variant-invariant while
$\diffop$ is the variant axis. The pair $(Q,\diffop)$ is, almost literally, a row
in Palace's integrator list. We pick up this thread, and the fold that consumes
the list, in \cref{ch:assembly}.
\end{remark}

A fourth operator, $\diffop=\Div$, completes the family on grounds of symmetry:
it is the div--div term, well-formed on the $\Hdiv$ space, and it fills the last
slot of the de Rham ladder $\Hone\xrightarrow{\Grad}\Hcurl\xrightarrow{\Curl}
\Hdiv\xrightarrow{\Div}\Ltwo$ that organizes \cref{ch:derham}. Among the five
analyses of this book it happens that no driver assembles a div--div stiffness
term, so it appears here as the symmetric completion of the family rather than as
a witnessed workhorse---but the template is the more compelling for being closed
under the full ladder of differential operators.

\section{The weak form as virtual work}
\label{sec:virtual-work}

We close the conceptual arc by returning to the energy principle promised in
\cref{sec:virtual-work}'s forward reference at the start of the chapter. The
bilinear-form family is not just a computational convenience; when $a$ is
symmetric and positive---as the electrostatic and magnetostatic stiffness forms
are---the weak problem $a(u,v)=\ell(v)$ is exactly the condition for a certain
energy to be minimized.

Define the energy functional
\begin{equation}
  J(w) \;=\; \tfrac12\,a(w,w) \;-\; \ell(w).
  \label{eq:energy}
\end{equation}
For the electrostatic problem, $\tfrac12 a(w,w)=\tfrac12\int_\dom
\varepsilon\abs{\Grad w}^2\dV$ is the electrostatic field energy stored by the
trial potential $w$, and $\ell(w)$ is the work done by the sources; $J$ is the
total potential energy. The claim is that the solution $u$ of the weak problem is
precisely the minimizer of $J$.

\begin{proposition}[Variational equivalence]
\label{prop:varequiv}
Let $a$ be a symmetric bilinear form ($a(w,z)=a(z,w)$) that is positive
($a(w,w)>0$ for $w\neq0$), and let $\ell$ be linear. Then $u$ solves the weak
problem $a(u,v)=\ell(v)$ for all $v$ \emph{if and only if} $u$ minimizes
$J(w)=\tfrac12 a(w,w)-\ell(w)$.
\end{proposition}

\begin{proof}
Perturb the candidate $u$ by an arbitrary admissible variation $v$, scaled by a
real parameter $t$, and expand $J(u+tv)$ using bilinearity and symmetry of $a$:
\[
  J(u+tv) = \tfrac12 a(u,u) + t\,a(u,v) + \tfrac{t^2}{2}a(v,v) - \ell(u) - t\,\ell(v).
\]
Group by powers of $t$:
\[
  J(u+tv) = J(u) + t\bigl(a(u,v)-\ell(v)\bigr) + \tfrac{t^2}{2}a(v,v).
\]
The coefficient of the \emph{linear} term in $t$ is exactly the weak-form
residual $a(u,v)-\ell(v)$. The quadratic coefficient $\tfrac12 a(v,v)$ is
positive (for $v\neq0$) by hypothesis, so as a function of $t$, $J(u+tv)$ is an
upward parabola. Its vertex sits at $t=0$---i.e.\ $u$ is the minimizer along the
direction $v$---\emph{if and only if} the linear coefficient vanishes,
$a(u,v)=\ell(v)$. Since $v$ was arbitrary, $u$ minimizes $J$ over all directions
exactly when it solves the weak problem.
\end{proof}

\begin{figure}[t]
\centering
\begin{tikzpicture}[scale=1.0]
  \draw[->,simGray] (-2.6,0) -- (2.8,0) node[right,font=\scriptsize]{$t$ (amount of perturbation $v$)};
  \draw[->,simGray] (0,-0.3) -- (0,2.7) node[above,font=\scriptsize]{$J(u+tv)$};
  % the parabola, vertex at t=0
  \draw[simBlue,very thick,domain=-2.3:2.3,samples=80]
    plot (\x,{0.42*\x*\x+0.35});
  % vertex marker
  \fill[simRust] (0,0.35) circle (2pt);
  \node[simRust,font=\small,below right] at (0,0.35) {$u$ (minimizer)};
  % tangent line at vertex (horizontal) = zero linear term
  \draw[simGreen,thick,densely dashed] (-1.6,0.35) -- (1.6,0.35);
  \node[simGreen,font=\scriptsize] at (1.7,0.2) {slope $=a(u,v)-\ell(v)=0$};
  % quadratic-term annotation
  \node[simBlue,font=\scriptsize] at (1.9,2.2) {curvature $=a(v,v)>0$};
\end{tikzpicture}
\caption[The weak form as energy minimization]{The energy
$J(u+tv)=J(u)+t\,(a(u,v)-\ell(v))+\tfrac{t^2}{2}a(v,v)$ along any perturbation
direction $v$ is an upward parabola (positive curvature $a(v,v)>0$). Its vertex
falls at $t=0$---meaning $u$ is the energy minimizer---precisely when the linear
slope $a(u,v)-\ell(v)$ vanishes, which is the weak equation. ``Solve the weak
form'' and ``minimize the energy against every virtual displacement'' are the
same statement; the weak form is the principle of virtual work.}
\label{fig:virtualwork}
\end{figure}

\Cref{prop:varequiv} is the rigorous form of the slogan we opened with. The
linear term $a(u,v)-\ell(v)$ is the first-order change in energy under a
``virtual'' perturbation $v$; demanding that it vanish for every $v$ is the
\emph{principle of virtual work}, and it is identically the weak form
(\cref{fig:virtualwork}). When $a$ is symmetric and positive, solving the weak
form is minimizing an energy---the form in which a physicist would have posed the
problem in the first place. (The driven and eigenmode problems are not energy
minimizations in this simple sense---their forms are indefinite or
complex---but the weak form survives intact as the master formulation; it is only
the energy \emph{interpretation} that is special to the symmetric-positive case.)

\section{From weak form to finite elements: the road ahead}

We have, in this chapter, taken five strong-form PDEs and reduced each to the
same abstract shape: \emph{find $u$ such that $a(u,v)=\ell(v)$ for all admissible
$v$}, with $a$ drawn from the one-parameter family $a(u,v)=(Q\,\diffop u,\diffop
v)$. Two large questions remain, and they organize the next several chapters.

First, what is the space of ``admissible'' functions, and how do we replace it
with a finite-dimensional one we can compute in? The weak form demands only one
derivative, which is exactly the regularity the piecewise-polynomial
finite-element spaces provide; \emph{which} space ($\Hone$, $\Hcurl$, $\Hdiv$)
goes with which differential operator is the content of \cref{ch:galerkin} and
\cref{ch:functionspaces}, and the way the four spaces interlock is the de Rham
complex of \cref{ch:derham}. Second, how does the bilinear form become a matrix?
Restricting $a$ to a finite-dimensional space spanned by basis functions turns
$a(u,v)=\ell(v)$ into a linear system $\mat{K}\vx=\vb$, assembled element by
element as a fold over the weak-form terms---the Galerkin method of
\cref{ch:galerkin} and the assembly machinery of \cref{ch:assembly}. The
essential/natural boundary distinction we drew here is implemented, in that
discrete setting, by the degree-of-freedom elimination of \cref{ch:bc}.

\summarysection
The strong form of a PDE asks too much: it demands a twice-differentiable
solution and enforces the equation pointwise. Multiplying by a test function $v$,
integrating over the domain, and integrating by parts trades one derivative on
the unknown $u$ for one on $v$, halving the smoothness demand and reorganizing
the equation into a \emph{bilinear form} $a(u,v)$ equal to a \emph{linear
functional} $\ell(v)$. In one dimension the integration by parts is the ordinary
formula and produces $\int u'v'=\int fv$; in three dimensions the divergence
theorem (Green's first identity) produces $a(u,v)=\int_\dom\varepsilon\Grad
u\cdot\Grad v$ plus a boundary integral; for vector fields the curl
integration-by-parts identity produces $a(\vu,\vv)=\int_\dom\nu\,\Curl\vu\cdot
\Curl\vv$, paired in eigenmode problems with the mass term
$\int_\dom\varepsilon\,\vu\cdot\vv$. The boundary integral splits into an
\emph{essential} (Dirichlet) part---enforced by constraining the test/trial
space, leaving no trace in the equation---and a \emph{natural} (Neumann) part,
which survives as a source in $\ell$. All of these are instances of one template,
$a(u,v)=(Q\,\diffop u,\diffop v)_{\Ltwo(\dom)}$, with $\diffop\in\{\text{value},
\Grad,\Curl,\Div\}$ and $Q$ a material coefficient---the four-term family that
every Palace driver assembles. When $a$ is symmetric and positive, the weak form
is equivalent to minimizing the energy $J=\tfrac12 a(u,u)-\ell(u)$: the principle
of virtual work.

\furtherreading
The classical reference for the variational formulation and its analysis is
Brenner and Scott~\citep{brennerscott2008}, which develops weak forms, Sobolev
regularity, and the Lax--Milgram and C\'ea theorems we have only gestured at.
For the curl--curl and vector cases specific to electromagnetics---Green's
identities for the curl operator, the $\Hcurl$ and $\Hdiv$ spaces, and the
boundary conditions of Maxwell problems---Monk~\citep{monk2003} is the standard
graduate treatment. Jin~\citep{jin2014} gives the same material from a
practicing computational-electromagnetics viewpoint, with the bilinear forms
written exactly as the integrators that appear in a solver. Readers wanting the
energy/virtual-work perspective in its mechanical home will find it in any
finite-element text rooted in structural analysis; the equivalence of
\cref{prop:varequiv} is there called the principle of minimum potential energy.

\exercisessection
\begin{enumerate}[nosep]
  \item \textbf{(1-D, from scratch.)} Derive the weak form of $-u''+u=f$ on
  $(0,1)$ with $u(0)=u(1)=0$. Identify the bilinear form $a(u,v)$ and the
  functional $\ell(v)$. How does $a$ differ from the pure-Poisson case
  \cref{eq:1d-bilinear}, and which two members of the family
  \cref{eq:family} does it combine?
  \item \textbf{(Boundary bookkeeping.)} Redo the 1-D derivation of
  \cref{sec:1d} but with the \emph{inhomogeneous} Dirichlet condition $u(0)=a$,
  $u(1)=b$. Show that the bilinear form is unchanged and explain where the data
  $a,b$ must enter instead. (This is the ``lifting'' of \cref{ch:bc}, in
  miniature.)
  \item \textbf{(Why $v$ vanishes.)} In the 1-D Neumann derivation
  \cref{eq:1d-neumann-weak} we kept $v(1)$ free but required $v(0)=0$. What weak
  form results if you mistakenly require $v(1)=0$ as well? Show that the
  prescribed flux $g$ then disappears from the equation entirely, and explain in
  one sentence why that makes the discrete problem wrong.
  \item \textbf{(Green's identity by hand.)} Starting from the product rule
  \cref{eq:prodrule}, derive Green's \emph{second} identity
  $\int_\dom (u\Lap v - v\Lap u)\dV = \oint_\bdry(u\,\partial_n v - v\,\partial_n
  u)\,\mathrm{d}S$ by applying Green's first identity twice (with the roles of
  $u$ and $v$ swapped) and subtracting. Here $\partial_n=\vn\cdot\Grad$.
  \item \textbf{(Reading the template.)} For each of the following bilinear
  forms, identify the coefficient $Q$, the differential operator $\diffop$, and
  the finite-element space on which it is well-formed: (a)
  $\int_\dom \mu^{-1}\Curl\vu\cdot\Curl\vv$; (b) $\int_\dom \rho\,u\,v$; (c)
  $\int_\dom \kappa\,\Grad T\cdot\Grad S$ (steady heat conduction). Which member
  of \cref{fig:family} is each?
  \item \textbf{(Symmetry and the energy.)} Show directly from the definition
  \cref{eq:family} that $a(u,v)=a(v,u)$ whenever the coefficient $Q$ is a
  symmetric matrix (or a scalar). Then exhibit a $2\times2$ matrix coefficient
  $Q$ for which $a(u,v)\neq a(v,u)$, and explain why \cref{prop:varequiv}'s
  energy interpretation fails for that $Q$.
  \item \textbf{(Eigenmode pencil.)} The eigenmode weak form is
  $a(\vu,\vv)=\omega^2 m(\vu,\vv)$ with $a$ the curl--curl form and $m$ the mass
  form \cref{eq:mass}. Writing both as members of the family
  \cref{eq:family}, give the pair $(Q,\diffop)$ for each. Discretized, this
  becomes $\mat{K}\vx=\lambda\mat{M}\vx$ with $\lambda=\omega^2$; which form
  produces $\mat{K}$ and which produces $\mat{M}$?
  \item \textbf{(Virtual work numerically.)} For the 1-D problem of
  \cref{wex:1dfull} (with exact solution $u=\sin\pi x$), compute the energy
  $J(u)=\tfrac12 a(u,u)-\ell(u)$ at the true solution. Then compute $J(w)$ for
  the perturbed trial function $w(x)=\sin\pi x + \tfrac1{10}\sin 2\pi x$ and
  confirm $J(w)>J(u)$, consistent with \cref{prop:varequiv}. (You may use
  $\int_0^1\sin m\pi x\,\sin n\pi x\dx=\tfrac12\delta_{mn}$ and the matching
  cosine identity.)
\end{enumerate}
